\section{Introduction}
\label{sec:intro}

When an AI alignment intervention fails to produce a measurable effect, what does that
mean? A team testing whether activation steering can elicit harmful outputs observes
$p > 0.05$ and concludes the model is robust. A safety evaluation shows no statistically
significant difference between a fine-tuned and base model and reports a clean bill
of health. These conclusions are drawn hundreds of times per year across the alignment
literature---and they may be wrong in the majority of cases.

The core problem is that classical null hypothesis significance testing (NHST) is
epistemically asymmetric: failing to reject the null tells us only ``we could not
detect an effect,'' not that the effect is absent. For alignment research, this
asymmetry has practical consequences. Teams stop investigating a potential vulnerability
after an inconclusive null test. Papers claim ``no jailbreak succeeded'' after testing
one class of prompts with $n=15$ trials. Safety evaluations declare robustness based
on benchmarks that \citet{kim2026fivnines} show are systematically saturated. Each
of these decisions treats experimental weakness as evidence of genuine safety.

{\bf Why this matters.} AI alignment research is accumulating a literature of null
results, and those null results are actively shaping deployment decisions, capability
evaluations, and research priorities. If the majority of those null results are
inconclusive rather than informative, the field is building on a shaky empirical
foundation. The inability to distinguish ``this model is robust'' from ``our experiment
was too weak to detect the vulnerability'' is not merely a statistical nuance---it is
a potential safety risk.

{\bf What existing work provides.} Bayesian methods for null hypothesis testing are
well-developed \citep{ly2021bayes,bartos2026efficient,schad2024accurate} and
equivalence testing frameworks provide principled tools for confirming
nulls in appropriate settings \citep{dale2024confirming}. The alignment literature
has begun addressing measurement challenges: \citet{kim2026fivnines} propose
importance sampling to estimate rare failure rates in saturated benchmarks, and
\citet{santos2026limits} formally prove that behavioral evaluation cannot uniquely
identify latent alignment under evaluation-aware agents. Sequential Bayesian
experimental design provides adaptive stopping rules for data collection
\citep{foster2025stopping,robbins2025robust}.

{\bf What is missing.} No existing framework connects these threads for alignment null
results specifically. Alignment papers rarely report effect sizes or retrospective
power. Intervention strength---the steering coefficient $\alpha$ in activation
engineering, the KL penalty in \rlhf, the specificity of a jailbreak prompt---is
treated as a binary ``we applied the intervention'' rather than as a continuous
parameter with implications for the expected effect size. Measurement sensitivity
is assumed rather than modeled. The result is that researchers lack a practical tool
to ask: ``given what we observed and how we designed this experiment, is this null
result evidence of robustness or evidence of an underpowered test?''

{\bf Our approach.} We develop \ours, a Bayesian framework that addresses this gap
directly. \ours formalizes two competing hypotheses for any alignment null result:
$\Hnull$ (true robustness: the effect size $\delta \approx 0$) and $\Halt$ (experimental
weakness: $\delta > 0$ but undetected due to insufficient power). For $\Halt$ we place
a calibrated Gaussian prior $\delta \sim \gN(0.3, 0.2^2)$ based on effect sizes
reported in the activation steering and \rlhf literatures. We derive a closed-form
Bayes factor $\mathrm{BF}_{01}$ exploiting Gaussian conjugacy, enabling rapid
retrospective analysis without MCMC. We then provide a $9 \times 9$ sensitivity
analysis over prior hyperparameters to verify that conclusions are not artifacts of
prior choice.

We validate \ours in three settings: \textbf{(1)} synthetic experiments where the
true effect size is known, confirming calibration under both $\Hnull$ and $\Halt$;
\textbf{(2)} retrospective analysis of 27 published alignment null results, testing
our central hypothesis that the majority are inconclusive; and \textbf{(3)}
prospective experiments with real \llama API calls under systematically varied
intervention strength, including a sequential Bayesian stopping analysis.

{\bf Main findings.} Applied to 27 published null results, \ours finds:
\begin{itemize}[leftmargin=*,itemsep=0pt,topsep=4pt]
    \item \mainresult of null results have $P(\text{weakness} \mid \text{data}) > 0.3$,
    far exceeding our pre-registered 40\% threshold.
    \item \emph{Zero} published null results constitute strong evidence of genuine
    robustness ($\mathrm{BF}_{01} > 10$); even the four best-supported studies only
    reach moderate evidence.
    \item Prospective \llama experiments reveal an ``inverse jailbreak'' effect:
    weak educational-framing interventions yield 20\% success while strong explicit
    jailbreak patterns yield 0\%---and \ours correctly classifies the strong-intervention
    null as inconclusive (not robust), preventing a false safety claim.
    \item Sequential analysis shows perfect agreement with post-hoc analysis
    ($\kappa = 1.000$), confirming the framework's prospective utility.
\end{itemize}

{\bf Contributions.} We make the following contributions:
\begin{itemize}[leftmargin=*,itemsep=0pt,topsep=4pt]
    \item We propose \ours, the first Bayesian framework specifically designed to
    diagnose whether alignment null results reflect genuine robustness or
    experimental insufficiency (\secref{sec:method}).
    \item We curate a dataset of 27 published alignment null results with extracted
    effect sizes and apply \ours retrospectively, finding systematic underpowering
    across the literature (\secref{sec:results}).
    \item We conduct prospective LLM experiments demonstrating the framework's
    practical value, including the discovery of an inverse jailbreak effect where
    stronger interventions produce null results that weaker interventions do not
    (\secref{sec:results}).
    \item We provide a prior sensitivity analysis showing that the ``inconclusive''
    classification is robust across a wide range of prior specifications, ruling out
    prior-sensitivity as an alternative explanation (\secref{sec:results}).
\end{itemize}
